\section{Empirical Evidence}
\label{sec:empirical}

We conducted three studies. Study~1 uses a two-condition suppression experiment across twelve models to test the ``last fingerprint'' hypothesis. Study~2 adds a third condition---explicit em dash prohibition---to test the dual-register mechanism directly. Study~3 compares a base model against its instruction-tuned counterpart to isolate the RLHF effect. A supplementary long-form analysis examines output length.

\subsection{Method}

We prompted instruction-tuned LLMs from five providers---Anthropic (Claude Opus~4.6, Claude Sonnet~4, Claude Haiku~3.5), OpenAI (GPT-4.1, GPT-4o, GPT-4o~Mini, GPT-5.4), Meta (Llama~3.1~8B~Instruct, Llama~3.3~70B~Instruct), Google (Gemini~2.5~Pro, Gemini~2.5~Flash), and DeepSeek (DeepSeek~V3)---with essay-writing tasks on ten everyday topics. Each model was run under two conditions:

\begin{itemize}[leftmargin=*]
\item \textbf{Unconstrained}: ``Write a 1000-word essay about [topic].''
\item \textbf{Prose-constrained}: Same, with the addition: ``Write in flowing prose paragraphs only. Do not use any markdown formatting, headers, bullet points, bold text, or lists.''
\end{itemize}

Each model received 10 topics per condition, yielding approximately 10,000 words per model per condition ($\sim$240,000 words total across all models and conditions).\footnote{Gemini~2.5~Pro was limited to three topics per condition ($\sim$3,100 words per condition) due to API rate limits; all other models completed the full ten topics.} For each sample, we counted Unicode em dash characters (U+2014) and overt markdown formatting features (headings, bullet points, bold text, numbered lists). All counts were normalized to frequency per 1,000 words.

All API calls used provider-default sampling parameters (temperature, top-$p$) with a maximum output length of 2,048 tokens. No frequency penalty, presence penalty, or system prompts were applied. For the base-vs-instruct comparison (Study~3), Llama~3.1~8B was run locally via Ollama using Q4\_K\_M quantization on identical hardware. All experiments were conducted between February and March 2026.

As a human baseline, we measured em dash frequency in eight published essays spanning literary criticism, journalism, and technical writing (57,232 words total). Human em dash usage varies substantially by genre and author: the weighted mean across our sample was 3.23 per 1,000 words (median 3.83; range 0.33--17.12).

\subsection{Results}

\begin{table}[ht]
\centering
\small
\begin{tabular}{llcccc}
\toprule
& & \multicolumn{2}{c}{\textbf{Em Dashes / 1K}} & \multicolumn{2}{c}{\textbf{MD Features / 1K}} \\
\cmidrule(lr){3-4} \cmidrule(lr){5-6}
\textbf{Model} & \textbf{Provider} & Unconstr. & Constr. & Unconstr. & Constr. \\
\midrule
GPT-4.1             & OpenAI    & 10.62 & 9.10 & 6.27 & 0.0 \\
Claude Opus 4.6     & Anthropic &  9.09 & 0.19 & 0.96 & 0.0 \\
Claude Sonnet 4     & Anthropic &  8.29 & 1.31 & 6.15 & 0.0 \\
Claude Haiku 3.5    & Anthropic &  7.51 & 0.18 & 5.36 & 0.9 \\
DeepSeek V3         & DeepSeek  &  6.95 & 5.41 & 1.47 & 0.0 \\
GPT-4o Mini         & OpenAI    &  4.16 & 4.23 & 6.03 & 0.0 \\
GPT-4o              & OpenAI    &  4.12 & 2.68 & 5.38 & 0.0 \\
Gemini 2.5 Pro      & Google    &  3.53 & 0.00 & 0.85 & 0.0 \\
GPT-5.4             & OpenAI    &  1.43 & 0.29 & 0.00 & 0.0 \\
Gemini 2.5 Flash    & Google    &  1.28 & 1.48 & 1.06 & 0.0 \\
Llama 3.1 8B Inst.  & Meta      &  0.00 & 0.00 & 1.91 & 0.0 \\
Llama 3.3 70B Inst. & Meta      &  0.00 & 0.00 & 0.00 & 0.0 \\
\midrule
Human baseline      &           & \multicolumn{2}{c}{3.23 (mean)} & \multicolumn{2}{c}{---} \\
\bottomrule
\end{tabular}
\caption{Two-condition suppression test across twelve instruction-tuned LLMs from five providers ($\sim$10,000 words per model per condition; $\sim$240,000 words total) and a human baseline (57,232 words across eight published essays). Markdown (MD) features include headings, bullet points, bold text, and numbered lists. Models are sorted by unconstrained em dash frequency.}
\label{tab:suppression}
\end{table}

Four patterns are visible in Table~\ref{tab:suppression}.

\textbf{Universal markdown suppression.} The prose-constrained instruction eliminates or nearly eliminates overt markdown features across all twelve models---from every provider, at every capability tier. Only Claude Haiku~3.5 retains a small residual (0.9 per 1,000 words); all other models reach zero. The suppression instruction works as intended.

\textbf{Differential em dash persistence.} Em dash behavior under suppression varies dramatically. GPT-4.1 drops only from 10.62 to 9.10 per 1,000 words (a 14\% reduction). DeepSeek~V3 drops from 6.95 to 5.41 (22\%). Anthropic models show the strongest suppression: Claude Opus~4.6 drops from 9.09 to 0.19 (98\%). Google's Gemini~2.5~Pro drops from 3.53 to 0.00---complete elimination. Meta's Llama models produce zero em dashes in either condition across approximately 40,000 words. We verified that Llama samples do not use alternative representations (double hyphens or en dashes); the models genuinely do not produce dash-mediated clause transitions. Two models with low baseline rates---GPT-4o~Mini (4.16 to 4.23) and Gemini~2.5~Flash (1.28 to 1.48)---show slight increases under suppression; at these frequencies the between-condition differences are within sampling variability and should not be interpreted as directional effects.

\textbf{Generational change within OpenAI.} GPT-4.1 produces 10.62 em dashes per 1,000 words unconstrained and 9.10 under suppression. GPT-5.4---OpenAI's most recent model---produces 1.43 unconstrained and 0.29 under suppression, with zero markdown features in either condition. This dramatic reduction across model generations is consistent with the mechanism being identified and addressed, as the paper's genealogy predicts: once named, the artifact becomes straightforward to tune.

\textbf{Em dash frequency varies within the human range.} The human baseline of 3.23 per 1,000 words (mean) with a range of 0.33--17.12 demonstrates that em dash frequency varies enormously across human writers and genres. Several LLMs produce em dashes at rates within or near this range. The finding is therefore not that LLMs uniformly ``overuse'' em dashes relative to all human writing, but rather that (a)~the em dash is the specific artifact that resists formatting suppression, (b)~em dash behavior varies systematically across providers in a pattern consistent with different RLHF calibrations, and (c)~some models produce em dashes at rates well above the human mean even under explicit suppression.

\subsection{Study 2: The Suppression Gradient}

The two-condition design shows that em dashes persist under markdown suppression. But a natural objection is that this merely reflects instruction specificity: the suppression instruction targets markdown features, not em dashes. To test this, we ran a three-condition experiment on four high-em-dash models, adding an explicit em dash suppression condition:

\begin{itemize}[leftmargin=*]
\item \textbf{Condition A}: Unconstrained.
\item \textbf{Condition B}: ``No markdown formatting, headers, bullet points, bold text, or lists.''
\item \textbf{Condition C}: Same as B, plus ``Do not use em dashes.''
\end{itemize}

\begin{table}[ht]
\centering
\small
\begin{tabular}{lccc}
\toprule
& \multicolumn{3}{c}{\textbf{Em Dashes / 1K words}} \\
\cmidrule(lr){2-4}
\textbf{Model} & Unconstr. & MD Suppr. & EM Suppr. \\
\midrule
GPT-4.1         & 11.51 & 8.20 & 3.86 \\
DeepSeek V3     &  8.66 & 4.75 & 1.57 \\
Claude Opus 4.6 &  8.46 & 0.19 & 0.00 \\
GPT-5.4         &  0.75 & 0.15 & 0.00 \\
\bottomrule
\end{tabular}
\caption{Three-condition suppression gradient. Condition~A: unconstrained. Condition~B: markdown suppression (no headers, bullets, bold, lists). Condition~C: markdown suppression plus explicit em dash prohibition. Markdown features were zero in conditions B and C for all models.}
\label{tab:gradient}
\end{table}

The results reveal a suppression gradient that directly tests the dual-register hypothesis. For models with high em dash production, the ``no markdown'' instruction (Condition~B) reduces but does not eliminate em dashes, because the instruction does not recognize the em dash as formatting. Explicitly naming the em dash (Condition~C) produces further reduction---but GPT-4.1 still retains 3.86 per 1,000 words and DeepSeek~V3 retains 1.57 even under direct prohibition. The structural impulse overrides the explicit instruction.

Claude Opus~4.6, by contrast, achieves complete suppression at Condition~B (0.19) and zero at Condition~C. GPT-5.4, already tuned to low baseline rates, reaches zero under either suppression condition. The gradient reveals that the em dash's resistance to suppression is not binary---it varies by model, and correlates with how deeply the RLHF procedure embedded the structural tendency.

\subsection{Study 3: Base versus Instruction-Tuned}

The genealogy claims that training data installs a latent structural tendency and RLHF amplifies it. To test this directly, we ran Llama~3.1~8B in both its base (pre-RLHF) and instruction-tuned configurations locally on identical hardware using the same quantization (Q4\_K\_M), with 10 essay topics at $\sim$1,000 words each.

\begin{table}[ht]
\centering
\small
\begin{tabular}{llccc}
\toprule
\textbf{Model} & \textbf{Type} & \textbf{Words} & \textbf{Em Dashes/1K} & \textbf{MD Features} \\
\midrule
Llama 3.1 8B & Base (pre-RLHF) & 8{,}186 & 0.49 & 28 \\
Llama 3.1 8B & Instruct (post-RLHF) & 9{,}241 & 0.00 & 5 \\
\bottomrule
\end{tabular}
\caption{Base versus instruction-tuned comparison. Same architecture, same weights, same quantization, same hardware. The base model produces a small number of em dashes (0.49/1K) and substantial markdown features (28). The instruct model produces zero em dashes and fewer markdown features.}
\label{tab:base}
\end{table}

The base model produces a small but nonzero number of em dashes (0.49 per 1,000 words---4 across 8,186 words) alongside substantial markdown formatting features (28 total). The instruction-tuned model produces zero em dashes and fewer markdown features (5 total). For the Llama family specifically, RLHF \textit{reduced} both artifacts rather than amplifying them.

This result has two implications. First, a latent tendency toward em dash production exists in the base model, consistent with the genealogy's claim that markdown-saturated training data installs a structural orientation. Second, the direction of the RLHF effect is provider-dependent: Meta's RLHF drove em dashes toward zero, while OpenAI's amplified them to 14.03 per 1,000 words (Table~\ref{tab:gradient}). The same base-level tendency can be amplified or suppressed depending on the fine-tuning procedure applied. RLHF is the variable that determines the outcome.

\subsection{Supplementary: Output Length}

We tested whether em dash frequency increases with output length by prompting four models with three essay topics at 5,000 words under all three suppression conditions. At this scale, the suppression gradient (Table~\ref{tab:gradient}) is amplified: GPT-4.1 reaches 14.03 per 1,000 words unconstrained and retains 6.97 even under explicit em dash prohibition across 9,331 words. Claude Opus~4.6 reaches 12.91 unconstrained but drops to zero under either suppression condition. The structural-joint mechanism predicts this amplification: longer outputs require more transitions, producing more points where the structural impulse can surface.

\subsection{Interpretation}

These measurements reframe the em dash pattern in three ways.

First, the em dash is confirmed as \textbf{the artifact that resists formatting suppression}. Across all five providers, the suppression instruction eliminates markdown features to zero. The em dash persists in models from Anthropic, OpenAI, and DeepSeek---at dramatically different rates---while Llama shows none and Gemini~2.5~Pro achieves complete suppression to zero. This is the predicted behavior of a dual-register artifact: the instruction to write prose catches structural formatting but passes over a mark that is already prose-legal.

Second, em dash frequency and suppression resistance function as a \textbf{signature of the RLHF procedure}. All twelve models were trained on markdown-saturated corpora, yet they produce starkly different em dash patterns. Sam Altman's acknowledgment that em dash frequency was deliberately tuned upward in ChatGPT~\cite{altman2024emdash} confirms that this lever exists. The Llama zero, the GPT-4.1-to-5.4 reduction, and the Anthropic suppression behavior all point to RLHF as the mechanism that determines the amplitude.

Third, the \textbf{cross-provider variation constitutes a taxonomy of fine-tuning approaches}. OpenAI's older models resist suppression (GPT-4.1 retains 9.10 under constraint); its newest model has been actively reduced (GPT-5.4 at 0.29). Anthropic models show high unconstrained rates but excellent suppression compliance. DeepSeek shows moderate suppression (6.95 to 5.41). Google and Meta show low or zero baseline rates. These differences are diagnostic of the engineering decisions each provider has made.
